\documentclass[12pt,a4paper]{article}

\usepackage[utf8]{inputenc}
\usepackage[T1]{fontenc}
\usepackage[french]{babel}
\usepackage{lmodern}
\usepackage[a4paper,margin=2.5cm]{geometry}
\usepackage{hyperref}
\usepackage{xcolor}
\usepackage{array}
\usepackage{booktabs}
\usepackage{tabularx}
\usepackage{enumitem}
\usepackage{amsmath,amssymb}
\usepackage{tcolorbox}
\usepackage{float}
\usepackage{parskip}
\usepackage{pgfplots}
\pgfplotsset{compat=1.18}

\definecolor{accent}{HTML}{2A5A8A}
\definecolor{muted}{HTML}{555555}
\definecolor{bgcode}{HTML}{F4F4F4}
\definecolor{warn}{HTML}{B85C00}
\definecolor{thbg}{HTML}{EEF3F8}
\definecolor{okbg}{HTML}{EEF9EE}
\definecolor{ok}{HTML}{2A7A2A}

\hypersetup{colorlinks=true, linkcolor=accent, urlcolor=accent, citecolor=accent}

\tcbuselibrary{skins,breakable}
\newtcolorbox{formulabox}{
  colback=thbg, colframe=accent, boxrule=0.8pt,
  left=8pt, right=8pt, top=6pt, bottom=6pt
}
\newtcolorbox{attention}{
  colback=okbg, colframe=ok, boxrule=0.8pt,
  left=8pt, right=8pt, top=6pt, bottom=6pt
}
\newtcolorbox{avertissement}{
  colback=yellow!8, colframe=warn, boxrule=0.8pt,
  left=8pt, right=8pt, top=6pt, bottom=6pt
}

\renewcommand{\arraystretch}{1.5}

\title{\textbf{Résultats et benchmark}}
\author{Samir El Khoumri, Aly Hachem-Reda, Sami Chbicheb}
\date{Mai 2026}

\begin{document}

\maketitle
\tableofcontents
\newpage

% ============================================================
\section{Résultats et benchmark}
% ============================================================

% ------------------------------------------------------------
\subsection{Présentation du projet de benchmark}
% ------------------------------------------------------------

\subsubsection{Choix du projet}

Pour évaluer RTS-LLM dans des conditions réalistes, nous avons retenu le
projet open source \textbf{spring-cucumber-rest-api}. Ce projet est une
API REST développée avec Spring Boot et intégralement testée par des
scénarios BDD écrits en Gherkin avec Cucumber. Il réunit plusieurs
caractéristiques qui en font un terrain d'évaluation pertinent :

\begin{itemize}[noitemsep]
  \item Les tests sont exclusivement des scénarios BDD fonctionnels,
        ce qui correspond exactement au contexte cible de RTS-LLM.
  \item Le code source est en Java, compatible avec l'analyse statique
        par JavaParser intégrée à l'outil.
  \item L'historique Git est suffisamment riche pour constituer un corpus
        d'évaluation représentatif.
  \item Le projet est public et reproductible, ce qui permet de vérifier
        indépendamment les résultats.
\end{itemize}

\subsubsection{Caractéristiques du corpus}

\begin{table}[H]
\centering
\begin{tabular}{ll}
\toprule
\textbf{Caractéristique}              & \textbf{Valeur} \\
\midrule
Nombre de scénarios BDD               & 12 \\
Nombre de commits dans l'historique   & 220 \\
Période couverte                      & Février 2017 -- Avril 2021 \\
Commits analysés (LLM seul)           & 30 \\
Commits analysés (hybride)            & 30 \\
Commits avec vérité terrain non vide  & 27 (les deux modes) \\
Commits retenus pour l'évaluation     & 6 (LLM seul) / 7 (hybride) \\
\bottomrule
\end{tabular}
\caption{Caractéristiques du corpus d'évaluation.}
\label{tab:corpus}
\end{table}

\noindent Les commits retenus sont ceux pour lesquels la vérité terrain
est non vide (\texttt{gt\_size > 0}), c'est-à-dire les commits qui
modifient effectivement du code traversé par au moins un scénario.
Les commits dont le diff ne touche que des fichiers non couverts
(documentation, configuration, tests de bas niveau) sont exclus car
ils ne permettent pas de mesurer la qualité de la sélection.

% ------------------------------------------------------------
\subsection{Configurations évaluées}
% ------------------------------------------------------------

Deux configurations de RTS-LLM ont été comparées, utilisant toutes les
deux le modèle \textbf{Mistral 7B-instruct (quantification q3\_K\_M)}
via Ollama en exécution locale sur GPU (NVIDIA RTX A500, 4 Go VRAM).
La version quantifiée (3,3 Go) a été choisie car la version complète
(4,4 Go) débordait sur le CPU, réduisant la vitesse d'inférence à
environ 5 tok/s contre 30 tok/s avec q3\_K\_M :

\begin{itemize}
  \item \textbf{LLM seul} : le diff et l'ensemble des scénarios sont
        soumis directement à Mistral, sans présélection statique.
        Le modèle décide seul quels scénarios retester.

  \item \textbf{Hybride (analyse statique + LLM)} : une analyse statique
        du call graph présélectionne les scénarios structurellement liés
        aux méthodes modifiées. Mistral reçoit uniquement ces candidats,
        accompagnés de la chaîne d'appel les reliant au code modifié,
        et joue un rôle d'élagueur.
\end{itemize}

Cette comparaison vise à mesurer l'apport concret de la couche statique
sur les trois métriques d'évaluation.

% ------------------------------------------------------------
\subsection{Résultats quantitatifs}
% ------------------------------------------------------------

\subsubsection{Anomalies détectées dans les métriques brutes}

\paragraph{Double-comptage dans la vérité terrain.}
Lors de l'analyse des résultats bruts, nous avons constaté que le
\texttt{gt\_size} moyen ($\approx 14$) dépassait le nombre réel de
scénarios du projet (12). Ce chiffre peut même atteindre~24 ou~30, ce
qui est mathématiquement impossible avec 12 scénarios.

L'origine est un bug dans le script de collecte de couverture :
\texttt{run\_per\_scenario\_coverage.sh} ajoutait à la fois
\texttt{src/test/resources} et \texttt{src/test} à sa liste de
répertoires de recherche. Comme le premier est un sous-répertoire du
second, le \texttt{grep} récursif trouvait chaque scénario \textbf{deux fois},
produisant 24 entrées dans l'index au lieu de~12. Chaque scénario
réellement impacté contribuait ainsi à 2~entrées dans~$G$, gonflant
\texttt{gt\_size} du facteur~2 et plafonnant mécaniquement le rappel à~0{,}5.

Le bug est corrigé dans la version actuelle du script. Comme la
re-collecte de couverture (une exécution JVM par scénario et par commit)
représente un coût calculatoire important, nous dérivons les métriques
corrigées analytiquement : le rappel corrigé vaut exactement deux fois
le rappel mesuré (puisque $G_{\text{réel}} = G_{\text{mesuré}} / 2$ et
que TP est inchangé). La précision n'est pas affectée.

\paragraph{Taux de NaN élevé (sélection vide).}
Sur les 27~commits avec vérité terrain non vide, 21~commits (LLM seul)
et 20~commits (hybride) produisent des métriques indéfinies (NaN).
Dans tous ces cas, \texttt{sel\_size = 0} : le modèle n'a sélectionné
aucun scénario. Un bug dans \texttt{eval\_history.sh} en est la cause
directe : lorsque la liste de sélection est vide, le script n'appelle
pas \texttt{GroundTruthBuilder} avec l'argument \texttt{--selection},
qui renvoie alors \texttt{fn = 0} au lieu de \texttt{fn = gt\_size},
provoquant une division par zéro ($R = 0/0 = \text{NaN}$). Ces commits
sont donc exclus des moyennes, ce qui surestime les performances réelles
de l'outil : les~cas de sélection vide correspondent à des rappels
nuls~($R = 0$) et constituent les pires scénarios.

\subsubsection{Tableau comparatif}

\begin{table}[H]
\centering
\begin{tabular}{lccccc}
\toprule
\textbf{Configuration} & \textbf{Commits valides} & \textbf{Précision} & \textbf{Rappel brut} & \textbf{Rappel corrigé} & \textbf{F1} \\
\midrule
LLM seul               & 6  & 0{,}650 & 0{,}110 & 0{,}220 & 0{,}168 \\
Hybride (statique + LLM) & 7 & \textbf{0{,}911} & \textbf{0{,}384} & \textbf{0{,}767} & \textbf{0{,}514} \\
\midrule
Gain absolu            & $+1$ & $+0{,}261$ & $+0{,}274$ & $+0{,}547$ & $+0{,}346$ \\
\bottomrule
\end{tabular}
\caption{Métriques moyennes sur les commits non-NaN avec \texttt{gt\_size > 0}
         (rappel corrigé $= $ rappel brut $\times 2$ pour corriger le double-comptage).
         Modèle : Mistral 7B-instruct q3\_K\_M.}
\label{tab:resultats}
\end{table}

\begin{avertissement}
Les 21/27 commits NaN (LLM seul) et 20/27 commits NaN (hybride) correspondent
tous à \texttt{sel\_size = 0} et représentent des rappels nuls ($R = 0$).
Les moyennes du tableau ci-dessus sont donc calculées \textbf{uniquement sur les
meilleurs commits} et surestiment les performances réelles.
Si l'on impute $R = 0, P = 0, F1 = 0$ aux commits NaN, le rappel corrigé moyen
global tombe à $0{,}049$ (LLM seul) et $0{,}199$ (hybride).
\end{avertissement}

\subsubsection{Distribution du rappel par commit}

Au-delà des moyennes, la distribution du rappel par commit révèle un
comportement plus nuancé (rappels corrigés, $\times 2$) :

\begin{table}[H]
\centering
\begin{tabular}{lcc}
\toprule
\textbf{Tranche de rappel corrigé} & \textbf{LLM seul} ($n=6$) & \textbf{Hybride} ($n=7$) \\
\midrule
$R = 0$                        & 1 commit  (17\,\%) & 0 commit   (0\,\%) \\
$R \in ]0;\, 0{,}5]$           & 4 commits (67\,\%) & 2 commits (29\,\%) \\
$R \in ]0{,}5;\, 1{,}0[$       & 1 commit  (17\,\%) & 1 commit  (14\,\%) \\
$R = 1$                        & 0 commit   (0\,\%) & 4 commits (57\,\%) \\
\bottomrule
\end{tabular}
\caption{Distribution du rappel corrigé par commit (commits non-NaN avec \texttt{gt\_size > 0}).}
\label{tab:distribution}
\end{table}

\noindent Le mode hybride atteint un rappel parfait ($R = 1$) sur 4 des 7 commits
valides (57\,\%), alors que le LLM seul n'y parvient jamais. Cette différence
s'explique par les chemins d'appel fournis par l'analyse statique : le LLM
dispose d'un contexte structurel précis qui lui permet de relier le diff à chaque
scénario impacté, au lieu de raisonner uniquement sur des ressemblances lexicales.
Les 4 commits avec $R = 1$ en mode hybride sont ceux qui modifient des méthodes
peu centrales, avec une vérité terrain réduite (2 scénarios) : l'analyse statique
isole correctement les candidats et le LLM les confirme tous.

% ------------------------------------------------------------
\subsection{Interprétation des résultats}
% ------------------------------------------------------------

\subsubsection{Apport de l'analyse statique}

L'analyse statique apporte un gain significatif sur toutes les métriques :
$+0{,}261$ en précision, $+0{,}547$ en rappel corrigé, et $+0{,}346$ en F1.
Ces gains s'expliquent par deux effets cumulatifs.

\textbf{Contexte structurel.} Le LLM reçoit la chaîne d'appel reliant
chaque candidat à la méthode modifiée. Il peut ainsi raisonner sur la
nature réelle de l'impact plutôt que sur des ressemblances lexicales dans
le diff, ce qui améliore à la fois la précision (0{,}650 $\to$ 0{,}911)
et le rappel.

\textbf{Réduction de la taille du prompt.} Sans analyse statique, le LLM
reçoit la totalité des 12 scénarios dans un prompt de 9\,000 à 11\,000
caractères. Avec le filtre statique, seuls les candidats structurellement
liés au diff sont transmis, réduisant la longueur du prompt de 60 à 80\,\%.
Cela améliore la qualité du raisonnement du modèle et réduit le taux de
sélection vide (de 21/27 à 20/27), bien que l'effet reste limité.

\textbf{Rappel parfait sur les commits peu centraux.} Le mode hybride
atteint $R = 1$ sur 4 commits dont la vérité terrain ne compte que
2 scénarios. L'analyse statique identifie exactement les scénarios
structurellement liés, et le LLM les confirme intégralement.

\subsubsection{Explication du faible rappel}

Le rappel corrigé de~0{,}220 (LLM seul, commits non-NaN) et le taux de NaN
de 78\,\% s'expliquent par trois causes complémentaires :

\begin{itemize}
  \item \textbf{Le conservatisme du LLM sur les grands prompts.}
        Sur 27~commits avec vérité terrain non vide, 21~renvoient une
        sélection vide, ce qui est interprété comme une absence totale
        de réponse utile. Le modèle Mistral~7B-instruct q3\_K\_M,
        confronté à un prompt de 9\,000 à 11\,000 caractères (12 scénarios
        + diff complet), tend à répondre avec un tableau JSON vide lorsqu'il
        ne détecte pas de lien \emph{direct} et évident entre le diff et un
        scénario.

  \item \textbf{La nature centrale du projet.} Avec seulement 12 scénarios
        BDD couvrant une API REST resserrée, un commit qui modifie une
        méthode centrale (gestion des requêtes HTTP, authentification,
        \texttt{AbstractBddStepDefinition}) peut structurellement impacter
        la quasi-totalité des scénarios. La vérité terrain atteint alors
        24 à 30 entrées (12 à 15 scénarios réels après correction), alors
        que l'outil en sélectionne 1 à 2.

  \item \textbf{Le bug NaN dans l'évaluation.}
        Lorsque \texttt{sel\_size = 0}, le script d'évaluation
        \texttt{eval\_history.sh} ne passe pas l'argument
        \texttt{--selection} à \texttt{GroundTruthBuilder}, ce qui
        produit \texttt{fn = 0} au lieu de \texttt{fn = gt\_size},
        et donc $R = 0/0 = \text{NaN}$. Ces commits sont exclus des
        moyennes, masquant les pires performances de l'outil.
\end{itemize}

\subsubsection{Analyse d'un cas concret}

Pour illustrer les situations où l'outil et la vérité terrain divergent,
prenons le commit \texttt{f1eea53} (\emph{``docs: Add javadoc + update
readme file''}). Ce commit supprime la classe
\texttt{BasicAuthAuthentication} et déplace sa méthode \texttt{authenticate}
dans \texttt{DefaultRestApiStepDefinitionTest}.

\begin{itemize}
  \item \textbf{Côté RTS-LLM} : l'analyse statique détecte la nouvelle
        méthode \texttt{authenticate} comme modifiée, remonte le call graph
        jusqu'aux scénarios d'authentification et en sélectionne 1 à 2.

  \item \textbf{Côté vérité terrain} : la couverture est collectée avant
        le commit, dans l'état du code parent. La nouvelle méthode n'existait
        pas encore à ce stade : aucun scénario ne peut l'avoir couverte.
        La vérité terrain est donc vide (\texttt{gt\_size = 0}), et le
        commit est exclu de l'évaluation.
\end{itemize}

Ce cas illustre parfaitement la limitation des méthodes ajoutées par le
diff : RTS-LLM a identifié un impact potentiellement réel (le code
d'authentification a changé), mais la méthode de mesure ne peut pas le
confirmer. Dans les métriques, cela se traduit par des faux positifs
apparents qui pénalisent l'outil.

% ------------------------------------------------------------
\subsection{Discussion et limites}
% ------------------------------------------------------------

\subsubsection{Limites du benchmark}

L'évaluation repose sur un seul projet. Cette contrainte est inhérente
aux difficultés de mise en place de la vérité terrain --- la collecte
de couverture par scénario nécessite d'exécuter l'intégralité de la
suite de tests pour chaque scénario individuellement, ce qui représente
un coût important en temps de calcul. Les résultats obtenus ne peuvent
donc pas être généralisés à d'autres projets sans campagne d'évaluation
complémentaire.

Par ailleurs, le projet retenu présente une particularité structurelle
forte : ses 12 scénarios couvrent massivement le même code applicatif
(une API REST de taille réduite). Cela produit des vérités terrain de
grande taille relative, qui mécaniquement font baisser le rappel moyen
de l'outil.

\subsubsection{Utilisabilité en production}

En considérant les 21~commits NaN comme des rappels nuls, le rappel
corrigé \emph{global} tombe à $0{,}049$ (LLM seul) et $0{,}199$ (hybride).
Ces chiffres sont insuffisants pour garantir la détection des régressions
en production : l'outil, dans son état actuel, ne peut pas remplacer une
suite complète de tests.

Le mode hybride, avec un rappel global de $0{,}199$ et une précision de
$0{,}911$ sur les commits où il sélectionne quelque chose, peut néanmoins
constituer un \textbf{complément utile} dans un pipeline CI : exécuté en
première passe rapide, il identifie avec une très haute précision les
scénarios les plus susceptibles d'être affectés par les commits qui
modifient des méthodes peu centrales. Pour les commits qui modifient du
code traversé par la majorité des scénarios, l'outil renvoie une
sélection vide et l'exécution complète de la suite reste nécessaire.

\subsubsection{Perspectives d'amélioration}

Plusieurs pistes permettraient d'améliorer les résultats :

\begin{itemize}[noitemsep]
  \item \textbf{Corriger le bug NaN dans l'évaluation} : lorsque
        \texttt{sel\_size = 0}, passer \texttt{--selection ""} à
        \texttt{GroundTruthBuilder} pour obtenir $R = 0$ au lieu de NaN.
        Cela donnerait une image fidèle des performances réelles et
        permettrait d'inclure les 21 commits actuellement exclus.

  \item \textbf{Améliorer la résolution du call graph} : traiter les
        fichiers supprimés par le diff et les méthodes nouvellement
        introduites, qui constituent les deux principales sources de
        faux négatifs statiques.

  \item \textbf{Affiner le prompt} : explorer des formulations qui
        incitent le LLM à être moins conservateur sur les grands prompts,
        ou découper les scénarios en lots pour éviter la saturation du
        contexte.

  \item \textbf{Étendre le benchmark} : évaluer l'outil sur plusieurs
        projets de natures différentes (tailles variées, domaines métier
        distincts, couplage scénario/code plus faible) pour obtenir des
        résultats généralisables.

  \item \textbf{Utiliser un modèle plus puissant} : Mistral 7B-instruct
        q3\_K\_M est un modèle compact contraint par le matériel disponible.
        Des modèles plus grands (13B, 34B) ou spécialisés sur le code
        pourraient réduire significativement le taux de sélection vide.

  \item \textbf{Re-collecter la couverture} avec le script corrigé pour
        obtenir un \texttt{gt\_size} réel (12 entrées au lieu de 24) et
        mesurer le rappel brut sans correction analytique.
\end{itemize}

\end{document}
